\documentclass[11pt,a4paper]{ctexart}
\usepackage[a4paper,margin=1.78cm,headheight=15pt]{geometry}
\usepackage{amsmath,amssymb,mathtools,booktabs,tabularx,array}
\usepackage[most]{tcolorbox}
\usepackage{tikz}
\usetikzlibrary{arrows.meta,positioning}
\usepackage{xcolor,fancyhdr,hyperref,bookmark,microtype}
\newcolumntype{L}[1]{>{\raggedright\arraybackslash}p{#1}}
\setlength{\parindent}{2em}\setlength{\parskip}{0.34em}\setlength{\emergencystretch}{3em}
\renewcommand{\arraystretch}{1.2}\setlength{\tabcolsep}{3pt}
\definecolor{deep}{RGB}{42,58,73}\definecolor{soft}{RGB}{245,247,249}
\definecolor{warn}{RGB}{136,68,63}
\colorlet{themebg}{white}\colorlet{themefg}{black}\colorlet{themegray}{gray!70}
\colorlet{themecurve}{blue!70!black}\colorlet{themealert}{red!70!black}\colorlet{themeamber}{orange!80!black}
\hypersetup{colorlinks=true,linkcolor=deep,pdftitle={H-B01 函数结构在运算中的传播}}
\pagestyle{fancy}\fancyhf{}
\lhead{\small\color{deep}\textbf{数学一 · H-B01}}
\rhead{\small\color{deep}结构传播}\cfoot{\small\thepage}
\renewcommand{\headrulewidth}{0.4pt}
\newtcolorbox{corebox}[1]{breakable,colback=soft,colframe=deep,title=\textbf{#1},boxrule=0.8pt,arc=2mm}
\newtcolorbox{warnbox}[1]{breakable,colback=white,colframe=warn,title=\textbf{#1},boxrule=0.7pt,arc=1.5mm}
\begin{document}
\begin{center}
{\LARGE\textbf{H-B01｜函数结构在运算中的传播}}\\[0.4em]
{\large 奇偶、周期、对称怎样穿过极限、微分、积分与展开}
\end{center}
\begin{quote}
专题01负责函数结构，专题02/03/05/11负责各自运算。本桥梁只回答：结构经过一次运算后，哪些保持、哪些翻转、哪些因常数、中心或定义域而失效？
\end{quote}
\section{母接口：结构 + 运算 + 资格}
\[
\boxed{\text{原结构}\to\text{运算作用}\to\text{中心/定义域/正则性审计}\to\text{新结构}}
\]
\begin{center}
\begin{tikzpicture}[node distance=.55cm and .55cm,
box/.style={draw=deep,rounded corners=1.5mm,minimum width=3.1cm,minimum height=.9cm,align=center,fill=soft,font=\small},
arr/.style={-{Latex[length=2mm]},thick,draw=deep}]
\node[box](a){函数对象\\奇偶/周期/对称};
\node[box,right=of a](b){运算\\极限/导数/积分/展开};
\node[box,right=of b](c){资格\\定义域/中心/常数};
\node[box,below=of c](d){结果结构\\保持/翻转/迁移};
\draw[arr](a)--(b);\draw[arr](b)--(c);\draw[arr](c)--(d);
\end{tikzpicture}
\end{center}
\section{四条稳定传播}
\begin{tabularx}{\linewidth}{L{0.25\linewidth}L{0.31\linewidth}X}
\toprule
运算 & 稳定结论 & 必查边界\\
\midrule
极限 & 对称定义域上，偶/奇结构可传给存在的极限 & 极限点、左右路径、定义域是否对称\\
求导 & 偶函数导数为奇，奇函数导数为偶；周期性保持 & 可导区间、周期端点、点值接缝\\
积分 & 奇函数对称区间积分为零；偶函数可折半 & 积分中心、原函数常数、区间是否对称\\
泰勒/傅里叶 & 展开中心决定对称基；奇偶性消掉一半系数 & 中心迁移、收敛域、端点/跳跃恢复\\
\bottomrule
\end{tabularx}
\subsection{求导传播：反射结构翻转，平移结构保持}
默认函数在所讨论区间满足相应可导性。若
\[
f(a+t)=f(a-t),
\]
对 $t$ 求导得到
\[
\boxed{f'(a+t)=-f'(a-t)},
\]
所以关于 $x=a$ 的轴对称经过一次求导变成关于 $(a,0)$ 的中心反对称；特别地，若 $f$ 在 $a$ 处可导，则
\[
\boxed{f'(a)=0}.
\]
这只是驻点条件，不自动保证 $a$ 是极值点，极值仍要交给专题04检查单调变化或更高阶信息。

若
\[
f(a+t)+f(a-t)=2b,
\]
求导后得到
\[
\boxed{f'(a+t)=f'(a-t)},
\]
所以一般中心对称经过一次求导变成关于 $x=a$ 的轴对称。于是原点奇偶只是特殊情形：偶函数导数为奇函数，奇函数导数为偶函数。

\begin{center}
\begin{tikzpicture}[color=themefg,text=themefg,>=Stealth]
  \tikzset{
    axis/.style={themegray,line width=0.7pt,-{Stealth[length=1.6mm]}},
    guide/.style={themegray,densely dashed,line width=0.7pt},
    maincurve/.style={themecurve,line width=1.8pt},
    shiftcurve/.style={themeamber,densely dashed,line width=1.2pt},
    op/.style={themeamber,line width=1.1pt,-{Stealth[length=1.9mm]}},
    backop/.style={themegray,line width=0.9pt,-{Stealth[length=1.7mm]}},
    tag/.style={fill=themebg,inner sep=1.5pt,text=themefg,font=\small},
    rowhead/.style={fill=themebg,inner sep=2pt,text=themefg,font=\small\bfseries}
  }

  % ═════════════════════════════════════════════════════════════════════════════
  % Row 1: Axis symmetric -> Derivative is Central anti-symmetric (shift y = 6.4)
  % ═════════════════════════════════════════════════════════════════════════════
  \begin{scope}[shift={(0,6.4)}]
    \node[rowhead,anchor=south] at (0,2.7) {原函数 $f$ 轴对称 $x=a$};
    \draw[axis] (-2.5,0) -- (2.7,0) node[below,font=\small] {$x$};
    \draw[axis] (0,-0.3) -- (0,2.6);
    \draw[guide] (0,-0.2) -- (0,2.4);
    \draw[maincurve] plot[domain=-2.15:2.15,samples=90] (\x,{0.35+0.32*\x*\x});
    \draw[shiftcurve] plot[domain=-2.0:2.0,samples=90] (\x,{0.95+0.32*\x*\x});
    \fill[themealert] (0,0.35) circle (1.6pt);
    \fill[themeamber] (0,0.95) circle (1.6pt);
    \node[tag,anchor=north] at (0,-0.05) {$x=a$};
    \node[tag,text=themecurve,anchor=south west] at (1.2,1.35) {$f$};
    \node[tag,text=themeamber,anchor=south west] at (0.9,2.05) {$f+C$};
    \node[tag,text=themegray,anchor=north,font=\footnotesize] at (0,-0.55) {竖直平移仍保对称轴 $x=a$};
  \end{scope}

  \begin{scope}[shift={(9.2,6.4)}]
    \node[rowhead,anchor=south] at (0,2.7) {导函数 $f'$ 中心反对称 $(a,0)$};
    \draw[axis] (-2.5,0) -- (2.7,0) node[below,font=\small] {$x$};
    \draw[axis] (0,-1.7) -- (0,1.85);
    \draw[guide] (0,-1.5) -- (0,1.65);
    \draw[maincurve] plot[domain=-2.0:2.0,samples=60] (\x,{0.72*\x});
    \node[tag,anchor=north] at (0,-0.05) {$x=a$};
    \node[tag,text=themecurve,anchor=south west] at (1.35,1.05) {$f'$};
    \fill[themealert] (0,0) circle (1.8pt);
    \node[tag,text=themegray,anchor=north,font=\footnotesize] at (0,-1.85) {中心固定在 $(a,0)$};
  \end{scope}

  \draw[op] (3.05,7.35) -- node[tag,above]{求导} (6.15,7.35);
  \draw[backop] (6.15,6.75) -- node[tag,below]{从 $a$ 标准化积分} (3.05,6.75);

  % ═════════════════════════════════════════════════════════════════════════════
  % Row 2: Central anti-symmetric -> Derivative is Axis symmetric (shift y = 0)
  % ═════════════════════════════════════════════════════════════════════════════
  \begin{scope}[shift={(0,0)}]
    \node[rowhead,anchor=south] at (0,2.5) {原函数 $f$ 中心反对称 $(a,0)$};
    \draw[axis] (-2.5,0) -- (2.7,0) node[below,font=\small] {$x$};
    \draw[axis] (0,-1.7) -- (0,2.35);
    \draw[guide] (0,-1.5) -- (0,2.15);
    \draw[maincurve] plot[domain=-1.9:1.9,samples=90] (\x,{0.14*\x*\x*\x+0.32*\x});
    \draw[shiftcurve] plot[domain=-1.9:1.9,samples=90] (\x,{0.75+0.14*\x*\x*\x+0.32*\x});
    \fill[themealert] (0,0) circle (1.8pt);
    \fill[themeamber] (0,0.75) circle (1.8pt);
    \node[tag,anchor=north] at (0,-0.05) {$x=a$};
    \node[tag,text=themeamber,anchor=west,font=\footnotesize] at (0.12,0.75) {$(a,C)$};
    \node[tag,text=themecurve,anchor=south west] at (1.1,0.95) {$f$};
    \node[tag,text=themeamber,anchor=south west] at (1.0,1.95) {$f+C$};
    \node[tag,text=themegray,anchor=north,font=\footnotesize] at (0,-1.85) {$+C$ 搬走中心至 $(a,C)$};
  \end{scope}

  \begin{scope}[shift={(9.2,0)}]
    \node[rowhead,anchor=south] at (0,2.5) {导函数 $f'$ 轴对称 $x=a$};
    \draw[axis] (-2.5,0) -- (2.7,0) node[below,font=\small] {$x$};
    \draw[axis] (0,-0.3) -- (0,2.4);
    \draw[guide] (0,-0.2) -- (0,2.2);
    \draw[maincurve] plot[domain=-2.05:2.05,samples=90] (\x,{0.35+0.48*\x*\x});
    \node[tag,anchor=north] at (0,-0.05) {$x=a$};
    \node[tag,text=themecurve,anchor=south west] at (0.8,1.65) {$f'$};
    \node[tag,text=themegray,anchor=north,font=\footnotesize] at (0,-0.55) {对称轴 $x=a$};
  \end{scope}

  \draw[op] (3.05,0.95) -- node[tag,above]{求导} (6.15,0.95);
  \draw[backop] (6.15,0.35) -- node[tag,below]{从 $a$ 标准化积分} (3.05,0.35);

  \node[font=\small,text=themegray,anchor=north] at (4.6,-2.45)
    {核心机制：$+C$ 仅竖直平移，故保轴对称（$x=a$ 不变），但把中心对称点从 $(a,0)$ 搬移至 $(a,C)$。};
\end{tikzpicture}

\end{center}

连续求导会让反射符号逐次翻转，因此若相应高阶导数存在：
\[
f\text{ 偶}\Longrightarrow f^{(2k+1)}(0)=0,
\qquad
f\text{ 奇}\Longrightarrow f^{(2k)}(0)=0.
\]
这也是后面 Taylor 系数被奇偶结构筛掉一半的局部来源。

平移结构则在求导下保持：若 $f(x+T)=f(x)$，则
\[
f'(x+T)=f'(x),
\]
但 $T$ 不必是 $f'$ 的最小正周期；若 $f(x+L)=-f(x)$，则同样有
\[
f'(x+L)=-f'(x).
\]
反方向必须处理积分常数：$f'$ 周期只能先推出 $f(x+T)-f(x)$ 为常数，不能直接推出 $f$ 周期。

还有一个常用的跨专题接口：若 $f$ 在每个长度为 $T$ 的周期区间上连续、区间内部可导且 $f(a+T)=f(a)$，则 Rolle 定理保证每个 $(a,a+T)$ 内至少存在一个 $\xi$ 使
\[
f'(\xi)=0.
\]
这里 H-B01 只提供“周期性制造相等端点值”这一接口；Rolle 定理本身、资格与存在点构造由专题04拥有。

\subsection{周期结构进入极限与连续：先缩回一个周期，再做局部性质}
周期性本身只说明重复，不自动提供极限或有界性；一旦叠加极限/连续性，结论会显著加强。

若 $f$ 以 $T>0$ 为周期，并且存在有限极限
\[
\lim_{x\to+\infty}f(x)=L,
\]
则对任意固定 $x$ 都有 $f(x)=f(x+nT)$。令 $n\to\infty$，右端趋于 $L$，因此
\[
\boxed{f(x)\equiv L}.
\]
所以非恒定周期函数不可能在 $+\infty$（或 $-\infty$）拥有有限极限。这里真正的接口是：\textbf{周期把一个固定点复制到任意远处，极限要求所有足够远的复制都逼近同一个值。}

若 $f$ 连续且以 $T>0$ 为周期，则只需在一个紧区间 $[0,T]$ 上使用连续函数的最值定理：存在 $m,M$ 使
\[
m\le f(x)\le M,\qquad x\in[0,T].
\]
周期性再把这个界和取到最大/最小值的点复制到整个实轴，因此
\[
\boxed{\text{连续周期函数在 }\mathbb R\text{ 上有界，并能取得全局最大值与最小值}.}
\]
注意删除“连续”以后结论不成立；周期函数可以在一个周期内无界。删除“周期”以后，闭区间上的有界性也不能自动升级成全实轴有界。

\begin{corebox}{原函数的奇偶传播必须分成两种方向}
设定义域是关于 $0$ 对称的连通区间。
\begin{itemize}
  \item 若 $f$ 为\textbf{偶函数}，则归一化原函数
  \[
  F_0(x)=\int_0^x f(t)\,dt
  \]
  为奇函数；一般原函数 $F_0+C$ 只有在 $C=0$ 时仍为奇函数。因此这里确实要审计积分常数。
  \item 若 $f$ 为\textbf{奇函数}，则任意原函数 $F$ 都是偶函数。因为
  \[
  H(x)=F(x)-F(-x),\qquad H'(x)=f(x)+f(-x)=0,
  \]
  且 $H(0)=0$，故 $F(x)=F(-x)$。给偶函数加任意常数仍是偶函数，所以这一方向\textbf{不会}被积分常数破坏。
\end{itemize}
旧 Source 中这两种情形必须分开；“积分常数会破坏原函数奇偶”只能用于“偶函数的奇原函数”这一方向，不能泛化。
\end{corebox}

\begin{corebox}{1999 年经典判断：四个选项其实在考四种不同传播机制}
设 $f$ 连续，$F'=f$，在通常的关于原点对称连通定义域上判断：
\begin{enumerate}
\item 若 $f$ 为奇函数，则任意原函数 $F$ 都为偶函数：\textbf{正确}。这是上面的“奇 $\xrightarrow{\int}$ 偶”方向，积分常数不会破坏偶性；
\item 若 $f$ 为偶函数，则任意原函数 $F$ 都为奇函数：\textbf{错误}。例如 $f\equiv1$，原函数 $F=x+C$；只有 $C=0$ 时为奇函数；
\item 若 $f$ 为周期函数，则任意原函数 $F$ 都为周期函数：\textbf{错误}。仍取 $f\equiv1$，它有任意正周期，但 $F=x+C$ 不周期；
\item 若 $f$ 为单调增函数，则原函数 $F$ 必单调增：\textbf{错误}。取 $f(x)=x$，它严格递增，但 $F(x)=x^2/2+C$ 在整个实轴上先减后增，因为 $F'=f$ 的符号会过零。
\end{enumerate}

所以这道题的统一读法不是“背 A”，而是：
\[
\boxed{\text{原函数的形状由 }F'=f\text{ 决定；奇偶要审计常数，周期要审计平均值，单调要审计 }f\text{ 的符号}.}
\]
\end{corebox}

\begin{tcolorbox}[title=直觉锚点：积分常数只是把整张图上下平移]
加常数 $C$ 不改变图像的左右形状，只做竖直平移。于是关于 $y$ 轴的镜面对称会原样保留，所以偶性不会被常数破坏；但关于原点的中心对称会被一起抬到 $(0,C)$，因此一般不再关于原点对称。换句话说：
\[
\boxed{\text{竖直平移保轴对称，但会搬走中心对称的中心}.}
\]
这比死记“哪种会被常数破坏”更稳定：忘掉结论时，画一条偶函数和一条奇函数，把图像整体上移即可重建。
\end{tcolorbox}

\subsection{变上限积分：结构传播的是“基点 + 区间”}
令
\[
G(x)=\int_0^x f(t)\,dt.
\]
若 $f$ 偶，则 $G$ 奇；若 $f$ 奇，则 $G$ 偶。生成原因不是口诀，而是把 $x\mapsto -x$ 同时作用到积分上限和积分变量：
\[
G(-x)=\int_0^{-x}f(t)\,dt=-\int_0^x f(-u)\,du.
\]
一般固定下限要把常数显式拆出来。写成
\[
H(x)=\int_a^x f(t)\,dt=G(x)-G(a).
\]
若 $f$ 奇，则 $G$ 偶，所以对任意固定 $a$，$H$ 仍为偶函数；若 $f$ 偶，则 $G$ 奇，因此 $H$ 为奇函数当且仅当
\[
G(a)=\int_0^a f(t)\,dt=0.
\]
所以“基点不在对称中心”不是一律破坏奇偶，而是必须先检查它引入的常数是否会搬动相应对称中心。

更一般地，双变限对象
\[
H(x)=\int_{\alpha(x)}^{\beta(x)}f(t)\,dt
\]
的结构由三件事共同决定：$f$ 的变换规律、上下限在 $x\mapsto -x$ 或反射下怎样交换、以及交换上下限产生的负号。只看被积函数奇偶不足以决定 $H$ 的奇偶。若题目进一步要求对这类对象求导，则切换到 H-B03 的 Leibniz 公式与正则性资格；H-B01 只负责结构怎样经过积分与复合传播。

\subsection{把积分基点放在对称中心：反射结构会发生一次翻转}
设
\[
F(x)=\int_a^x f(t)\,dt.
\]
若 $f$ 关于 $x=a$ 轴对称，即
\[
f(a+t)=f(a-t),
\]
则
\[
\boxed{F(a+t)=-F(a-t)},
\]
所以标准化积分把“关于 $x=a$ 的轴对称”变成“关于 $(a,0)$ 的中心反对称”。

反之，若
\[
f(a+t)=-f(a-t),
\]
则
\[
\boxed{F(a+t)=F(a-t)},
\]
所以标准化积分把“关于 $(a,0)$ 的中心反对称”变成“关于 $x=a$ 的轴对称”。

这说明积分传播的中心不能省略：若下限不是对称中心，先拆成“以中心为下限的标准化积分 + 常数”，再判断常数是否搬动了对称中心。

\subsection{复合变上限：先积分生成中间函数，再研究复合}
对
\[
H(x)=\int_0^{\varphi(x)}f(t)\,dt,
\]
先定义
\[
F(u)=\int_0^u f(t)\,dt,
\qquad H=F\circ\varphi.
\]
于是可以把两层机制分开：
\[
\boxed{f\xrightarrow{\text{标准化积分}}F\xrightarrow{\text{复合 }\varphi}H}.
\]
因此在合法定义域上：
\begin{itemize}
\item $\varphi$ 为偶函数时，无论 $F$ 怎样，$H$ 都为偶函数；
\item $f$ 偶、$\varphi$ 奇时，$F$ 奇，从而 $H$ 奇；
\item $f$ 奇、$\varphi$ 奇时，$F$ 偶，从而 $H$ 偶。
\end{itemize}
这三条不是独立口诀，而是“积分传播一次 + 复合传播一次”的直接组合。

若
\[
\alpha(-x)=-\beta(x),\qquad \beta(-x)=-\alpha(x),
\]
则对
\[
H(x)=\int_{\alpha(x)}^{\beta(x)}f(t)\,dt
\]
作 $t=-u$ 后得到
\[
H(-x)=\int_{\alpha(x)}^{\beta(x)}f(-u)\,du.
\]
所以 $f$ 偶时 $H$ 偶，$f$ 奇时 $H$ 奇。若上下限本来就是 $[-\varphi(x),\varphi(x)]$，则奇被积函数直接抵消为零；偶被积函数则可折成两倍的单变上限积分。

\subsection{对称区间定积分：轴对称加倍，中心对称读平均高度}
若积分区间关于 $a$ 对称，则：
\[
f(a+t)=f(a-t)
\Longrightarrow
\boxed{\int_{a-r}^{a+r}f(x)\,dx=2\int_a^{a+r}f(x)\,dx};
\]
\[
f(a+t)=-f(a-t)
\Longrightarrow
\boxed{\int_{a-r}^{a+r}f(x)\,dx=0}.
\]
更一般地，若函数关于 $(a,b)$ 中心对称，即
\[
f(a+t)+f(a-t)=2b,
\]
则
\[
\boxed{\int_{a-r}^{a+r}f(x)\,dx=2br}.
\]
因此一般中心对称在对称区间上的平均函数值恰好等于中心高度 $b$。这条关系比只记“奇函数积分为零”更有生成力：把中心平移到 $(a,b)$ 后，奇偶积分只是它的特殊情形。

\section{周期、平均值与线性漂移}
若 $f$ 以 $T$ 为周期，则任何连续一周期积分都与起点无关：
\[
\int_x^{x+T}f(t)\,dt=\int_0^T f(t)\,dt=:M_T.
\]
于是变上限积分 $F(x)=\int_{x_0}^{x}f(t)\,dt$ 满足
\[
F(x+T)-F(x)=M_T.
\]
这给出比“原函数是否周期”更强的结构：
\[
\boxed{F(x)=\frac{M_T}{T}x+P(x)},\qquad P(x+T)=P(x).
\]
所以周期函数的原函数一般是“\textbf{线性漂移 + 周期波动}”；只有平均值 $M_T/T=0$ 时，原函数才可能本身周期。

\begin{center}
\begin{tikzpicture}[color=themefg,text=themefg,>=Stealth,x=0.95cm,y=1.0cm]
  \tikzset{axis/.style={themegray,line width=.7pt,-{Stealth[length=1.6mm]}}, curve/.style={themecurve,line width=2pt}, base/.style={themeamber,densely dashed,line width=1.1pt}, tag/.style={fill=themebg,inner sep=1pt,font=\small}}
  % left: zero mean
  \begin{scope}[shift={(0,0)}]
    \draw[axis] (-.4,0)--(7.2,0) node[below] {$x$};
    \draw[axis] (0,-2.0)--(0,2.25);
    \draw[curve] plot[domain=0:6.7,samples=140] (\x,{1.05*sin(120*\x)});
    \draw[themegray,densely dashed] (3, -1.75)--(3,1.75);
    \draw[themegray,densely dashed] (6, -1.75)--(6,1.75);
    \node[tag,anchor=north] at (3,0) {$T$}; \node[tag,anchor=north] at (6,0) {$2T$};
    \node[tag,anchor=south west] at (.6,1.35) {$M_T=0$};
    \node[font=\small,anchor=north] at (3.3,-2.28) {原函数只保留周期波动};
  \end{scope}
  % right: nonzero mean = drift + periodic
  \begin{scope}[shift={(9.2,0)}]
    \draw[axis] (-.4,0)--(7.2,0) node[below] {$x$};
    \draw[axis] (0,-.8)--(0,5.2);
    \draw[base] (0,.35)--(6.8,4.1);
    \draw[curve] plot[domain=0:6.7,samples=140] (\x,{.35+.55*\x+.72*sin(120*\x)});
    \draw[themegray,densely dashed] (3,0)--(3,4.2);
    \draw[themegray,densely dashed] (6,0)--(6,4.8);
    \node[tag,anchor=north] at (3,0) {$T$}; \node[tag,anchor=north] at (6,0) {$2T$};
    \node[tag,anchor=south west] at (.55,4.5) {$M_T\ne0$};
    \node[tag,anchor=south west,text=themeamber] at (4.1,3.15) {$\dfrac{M_T}{T}x$};
    \node[font=\small,anchor=north] at (3.3,-1.08) {$F(x)=\dfrac{M_T}{T}x+P(x)$};
  \end{scope}
  \node[font=\small,text=themegray,anchor=north] at (7.9,-3.1)
    {$F(x+T)-F(x)=M_T$：一周期平均值决定是否出现长期线性漂移。};
\end{tikzpicture}

\end{center}

有两类常见结构会自动把周期平均值压成零：
\begin{itemize}
\item 若 $f$ 同时是奇函数和 $T$-周期函数，则在 $[-T/2,T/2]$ 上由奇性抵消，而完整周期积分与起点无关，因此任意长度为 $T$ 的完整周期积分都为零；
\item 若 $f(x+L)=-f(x)$，则 $2L$ 是周期，前后两个半周期的积分互相抵消，因此任意长度为 $2L$ 的完整周期积分都为零。
\end{itemize}
两种情况下，相应原函数都不存在长期线性漂移，并可保持对应完整周期。

反方向也有稳定接口：若 $f'$ 以 $T$ 为周期，则
\[
\frac{d}{dx}\bigl[f(x+T)-f(x)\bigr]=0,
\]
故在合法连通区间上
\[
f(x+T)-f(x)=C,
\]
即 $f(x)=\frac CT x+P(x)$，其中 $P$ 为 $T$ 周期函数。这里的重点是“导数周期”只强制\textbf{增量恒定}，不自动强制原函数周期。

平移或反射后的函数应先变换中心，再谈奇偶和周期，不能把关于原点的口诀直接搬到新中心。

\subsection{多个反射条件：先复合几何变换，再读函数值关系}
专题01拥有反射/平移本体，本桥梁只保留其进入微积分运算时的接口：两个对称轴复合为平移；两个同高度对称中心复合为平移；若两个中心高度不同，则得到“横向平移 + 纵向漂移”，不是周期。轴对称与中心对称复合时常先得到关于某条基准线的翻转，再迭代一次得到真正周期。凡是要把这些结构带进积分或展开，先写出变换后的\textbf{函数值等式}，再做运算传播。
\section{展开系数：对称性不是装饰，而是系数过滤器}
若函数关于展开中心具有奇偶结构，则有限 Taylor 与 Fourier 系数会被系统性消去：偶函数在 $0$ 的奇阶导数为零，奇函数的偶阶导数为零；在对称区间上，偶函数的 Fourier 正弦系数为零，奇函数的余弦系数和常数项为零。

一般对称只需要先把中心搬到原点。若 $f$ 关于 $x=a$ 轴对称，则 $u\mapsto f(a+u)$ 为偶函数，因此在相应 Taylor 资格成立时只出现偶次幂：
\[
\boxed{f(x)=c_0+c_2(x-a)^2+c_4(x-a)^4+\cdots}.
\]
若 $f$ 关于 $(a,b)$ 中心对称，则 $g(u)=f(a+u)-b$ 为奇函数，因此
\[
\boxed{f(x)=b+c_1(x-a)+c_3(x-a)^3+\cdots}.
\]
所以“展开中心”不是装饰参数，而是把一般对称重新坐标化为原点奇偶的关键。

Fourier 中同样读取输入平移对基函数的作用。若 $f(x+L)=-f(x)$，把它看成 $2L$-周期函数时，平移 $L$ 会让第 $n$ 个 Fourier 模态获得因子 $(-1)^n$；反号结构只允许与 $-1$ 相容的奇数阶模态保留。这里 H-B01 只拥有“结构怎样过滤系数”的接口，具体 Fourier 系数计算、收敛与和函数恢复由专题11拥有。

这条接口的正确方向是
\[
\boxed{\text{函数结构}\to\text{基函数在该变换下的符号}\to\text{系数消失}},
\]
而不是“看到某些系数为零就无条件反推全局奇偶”。反推还需要展开确实恢复原函数、定义域关于中心对称等资格；这些收敛/恢复条件分别由专题11 与 H-B05 拥有。

\section{最小例子：周期能否传给原函数}
$f(x)=\sin x$ 的周期为 $2\pi$，且一周期积分为零，原函数 $F(x)=-\cos x$ 仍可取为周期函数。反之 $f(x)=1$ 也具有任意正周期，但
\[
F(x+T)-F(x)=\int_x^{x+T}1\,dt=T\ne0,
\]
所以任何原函数都不周期。题面把奇偶/周期带入积分、求导或展开时，先调用“结构 + 运算 + 资格”接口。

\section{反向边界与调用检查}
\begin{itemize}
\item 形式表达式对称不等于定义域对称；
\item 导数结构不等于原函数结构；
\item Taylor 的奇偶项消失不等于 Fourier 在每点恢复；
\item 周期函数的原函数不必周期。
\end{itemize}
遇到结构传播题，按“对象身份—运算—资格—不变量”复核，并在一个中心点或一个周期上做简单值检查。
\end{document}
